% !TeX root = ../papers/nc-p3.tex
\section{被約二十点を通る反転二窓の成分}
\label{sec:reversible-component}
$E,F$を四次元とし、$\delta$を命題\ref{prop:finite-flip-obstructions}の二重交代化とする。
第一関係空間$R\in\Gr(6,E\otimes F)$に対し、$\delta|_{\operatorname{Sym}^2R}$の階数が
20以下となる小行列式スキームを$Z$とする。以下の次元は以前のClifford点での
一次変形の次元を用いず、この$Z$の局所環から計算する。

\begin{theorem}[反転二窓の成分と幾何学的一般点]\label{thm:reversible-component}
被約二十点例の第一窓$R_*$を通る$Z$の成分は一つであり、その次元は44である。
$R_*$での接次元は45で、滑らかな周囲の45変数の形式座標における局所方程式は
階数6の二次式を初項とする一式である。
この成分の幾何学的一般点では、第一窓は四辺が有限な標準Koszul AS正則対象へ一意に延び、
周期は2を割り、$8\le\rho\le10$である。
第一窓の基底変換を除いた成分の有理商の次元は14で、その幾何学的一般点が表す添字保持同型類はこの商に従う。
\end{theorem}
\begin{proof}
例で$\delta$の核は可逆な六行対称行列$H$が張る直線である。
$R_*$の選んだ基底列を$P$ ($16\times6$)とし、$C=PHP^{\mathsf t}$と書く。
階数がちょうど20の近傍では、核直線は小行列式から正則に復元される。
従って$Z$はそこで、対称階数6の射影行列$[C]$で$\delta(C)=0$かつ上記の開条件を満たすものと
スキームとして同型である。逆射は$C$の像を取る。
列が独立な行列$P'$から$[P' H P'^{\mathsf t}]$への射は、相対次元16の滑らかな被覆である。
実際、像空間を局所的に自明化すれば可逆な対称形式を基底変換で$H$へ移す問題となり、
標数零ではエタール局所的に実行できる。ファイバーは直交相似群の捩れで、次元は$15+1$である。

この被覆の96変数を$P'=P+YH^{-1}$と取る。36個の方程式はちょうど
\[
 G(Y)+\delta(YH^{-1}Y^{\mathsf t})=0,\qquad
 G(Y)=\delta(PY^{\mathsf t}+YP^{\mathsf t})
\]
である。$G$の階数は35で、その余核の双対を張る行を$\lambda$とすると、
$\ker G$上の二次式 $q(Y)=\lambda\delta(YH^{-1}Y^{\mathsf t})$の階数は6である。
\path{computations/verify_reversible_component.py}では第一・第二座標の積の核を$P$、
二重交代化の核を$H$としてこの二つの階数を有理数で検算する。
形式陰関数定理で独立な35式を解くと、残る61変数の一式の初項は$q$である。
独立式の解の二次以上の補正に$G$を施しても$\lambda$で消えるためである。
階数6の斉次二次式は既約で、完備局所環の随伴次数環は多項式環をこの二次式で割った整域となる。
従って被覆の局所環は整域で、次元60、接次元61である。
滑らかな16変数を除けば元の局所環は次元44、接次元45で、一つの成分を持つ。
完備局所環への滑らかな変数の添加は最小表示の関係式数や二次初項の階数を変えないので、
元の局所方程式についても主張を得る。

対称行列$C$を$(0,2)\mid(1,3)$で読み替え、対称な添字対それぞれ十個に制限した行列を$B(C)$とする。
例では$B(C)$の階数は6である。上記検算では$Y\in\ker G$で$q(Y)=0$かつ、
$B(PY^{\mathsf t}+YP^{\mathsf t})$が$\ker B(C)$から$\operatorname{coker}B(C)$に
階数2以上を誘導する有理方向を求める。可逆な六行小行列を局所消去すれば、残る四行四列ブロックの
ある二次小行列式の二次初項がこの方向で非零である。
従って初項は$q$の倍数ではなく、局所環でもこの小行列式は非零である。
一般点の第三階数は8以上となり、上界は両端対称性から従う。

例の近傍で四窓の一因子階数は4、両隣接切断の階数は6である。
後者は対称性と$C$の階数6による。三次・四次交叉は上半連続で、
$w=C$とその縮約からそれぞれの下界4と1も成立するので、開近傍で例と同じ次元になる。
定理\ref{thm:frobenius-tensor-reconstruction}を適用すると、二次双対は
各次数の次元が一定の有限Frobenius代数族になる。実際、四頂点に折り返した各成分は
関係を挿入する行列の余核で、五次は零、関係が二次のため以後も零である。
階数一定の各成分と積はこの近傍でベクトル束とその射を定める。
四つの点対応は$R,R^{\rm op}$の反復であり、例での有限エタール性も開条件である。

内部次数$n$を固定した相対被約bar複体は有界なベクトル束複体である。
例ではKoszulなので全ての非対角ホモロジーは零であり、固定した$n$ごとに同じ消滅が
例を含む開集合で成立する。既約成分の幾何学的一般点はこれら全てに属するので、
そこでは全次数でKoszulである。有限Frobenius双対の完全対合から左右のAS条件と標準Hilbert級数が従う。
同じ証明は非可算な代数閉体上の非常に一般の閉点にも適用できる。
単一の開近傍の全閉点で全次数を認証したことにはならない。

両射影で定理\ref{thm:finite-all-square-roots}の剛性が成立する開集合も例を含む
（第二射影は例の二十行を組ごとに交換したもの）。
そこで各AS延長の次の窓は反転窓に限られ、二窓の一意性から第一窓の軌道が添字保持同型を決める。
両射影の点列は射影枠を含み、第一窓の安定化群は二十点への置換に忠実に作用する。
両射影の枠を点ごとに固定する射影変換は恒等だからである。
従って$\PGL(E)\times\PGL(F)$の軌道次元は30で、有理商の次元は$44-30=14$である。
\end{proof}

\begin{corollary}[対称部分に横断的な二次錐]\label{cor:reversible-quadratic-cone}
$Z$の$R_*$における完備局所環は、39個の自由変数と六変数の非退化二次錐の積である。
特異部分の形式近傍は、第一窓を対称化できる完全交叉の部分に一致する。
\end{corollary}
\begin{proof}
$E=F$と同定する。$R\subset\operatorname{Sym}^2 E$の開部分
（直交する四二次式に基点がないもの）は24次元で、対称完全交叉の定理から$Z$に入る。
二つの空間の相対的同定を$\PGL_4$で動かすと、例の近くで滑らかな39次元部分を得る。
実際、$(1\otimes g)R$が対称に入る条件は$g$の十六係数で線型で、$R_*$で階数15である
（前記検算）。階数15の開部分では唯一の対称化写像の直線を小行列式から復元でき、可逆性も開である。
ここで階数15の開部分は対称化可能な軌跡内で取る。
対称な各$R$では$G$の像は、$\bigwedge^2 E\otimes\bigwedge^2 E\to\bigwedge^4 E$
の積で消える。実際、$P$の各列は対称二階テンソルであり、一列と任意の列との積を四文字で
全交代化すると零である。従ってこの39次元部分に沿って微分の階数は35以下で、
近傍では35となる。ゆえに全体が特異軌跡に入る。

一式の二次初項の非退化な六変数を取り、残る39変数をパラメータにする。
六つの偏微分を形式陰関数定理で解いて臨界点を原点へ平行移動し、六変数を含む高次項を
次数順に非退化二次式の偏微分で消せば、方程式は$q_6(z)+g(t)$となる（標数零）。
特異部分は$z=0$に入るが、上述の39次元部分の完備局所環は滑らかで同次元なので、
その部分は$z=0$全体を占める。従って$g=0$であり、特異部分のスキーム構造も主張通りである。
\end{proof}

\begin{corollary}[一般最小周期]\label{cor:reversible-generic-period}
この成分の幾何学的一般点のAS延長の最小周期は2である。
\end{corollary}
\begin{proof}
一移動との添字保持同型があれば、第一窓は基底変換で$R^{\rm op}$に移る。
この条件は例の近傍で閉である。点を順序付けたエタール被覆で、各順列に対し、
二つの射影行列が第一・第二射影の対応する点を入れ替える条件を課す。
行列を射影化した全線型写像にコンパクト化し、行の比例を二次小行列式で表す。
この閉条件で特異な射影行列は現れない。実際、三次評価の独立性からベクトル次元
$r=1,2,3$の空間に属する行は高々$1,4,10$本である。
階数$r$の行列は次元$4-r$の核以外の全行を次元$r$の像に送るが、両方の上界の和は20未満である。
従って固有射の像は求める閉集合である。点対応から双一次式の空間も復元できるので、
全点を対応させることは窓を対応させることと同値である。

例の第一射影の二十点の自己同型群は$S_5$である。
七点以上を含む平面は先の十平面で尽くされることは、独立な三行が張る全平面を列挙して検査できる。
その法線の位置は五頂点の完全グラフの辺であり、三法線が従属なのは三角形の三辺の場合に限る。
従って置換は辺の接続関係を保ち、互いに接続する四辺の星から五頂点を復元できる。
全法線を固定する射影変換は、座標とその差を固定するので恒等である。
逆に$v_4$の符号を変えて五つの枠ベクトルの和を零にすれば、その任意の置換は
各対の和・差の二点とその相手を保つ。従って窓を保つ基底変換は対角な$S_5$であり、
反転と合成して窓を保つものはこれと因子交換との積で尽くされる。

射影スカラーを含めた基底変換軌道の接空間を$O$とする。$C$の対称行列表示ではこれは
$LC+CL^{\mathsf t}$、$L=u\otimes1+1\otimes v$が張る31次元空間である。
$S_5$の七共役類の各代表と因子交換との積を$J$とすると、検算ファイルでは$JCJ^{\mathsf t}=C$となる。
各代表に対し$Y\in\ker G$のある基底列で$D=YP^{\mathsf t}+PY^{\mathsf t}$とおくと、
$JDJ^{\mathsf t}-D$を$O$に加えた階数が標数101で32になる。
軌道次元と合わせ、標数零でも$JDJ^{\mathsf t}-D\in O$は真の線型条件である。
共役で他の全置換にも同じことがいえる。

階数6の既約二次錐$q=0$は$\ker G$を線型に張るので、これら有限個の真の部分空間を
全て避ける滑らかな方向を取れる。これは形式曲線の一階項になる。
実際、先に35式を消去した残りの式に変数$t(v+u)$（$v$は選んだ方向、$u(0)=0$）を代入し、
$t^2$で割った式を$q$の$v$での非零偏微分で解けばよい。
この曲線の一般点で反転窓に同型なら、二十点の順序を固定した後には、ある一定の順列で対応する。
特殊点でも行列が可逆であり、対応する五点の射影枠から変換は形式曲線上に一意に延びる。
従って射影スカラーを含めて一次で比較すると、特殊点の対応$J$に対し
$JDJ^{\mathsf t}-D\in O$となって矛盾する。
反転同型の閉集合は成分全体を含まず、一般点は一周期にならない。二周期性は命題\ref{prop:finite-flip-obstructions}による。
\end{proof}

\begin{proposition}[第三階数6と7の局所イデアル]\label{prop:reversible-low-opposite-ideals}
$Z$の$R_*$での局所環で、対称化可能な部分の被約イデアルを$I$とする。
第三階数$\rho\le6$の小行列式イデアルは$I$、$\rho\le7$のものは$I^2$である。
従ってこの点の開近傍では、第三階数6の点は対称部分に尽き、ちょうど7の点は存在しない。
\end{proposition}
\begin{proof}
$B(C)$の可逆な六行六列部分を消去し、残りを四行四列のSchur補行列$S$とする。
対称完全交叉では三つの二対二切断の階数が同じ6なので、$S$の全成分は$I$に入る。
$R_*$での$B(C)$は対称であり、$K$をその四次元核の基底行列とする。核と余核の基底を合わせると
$S$の微分は$T(Y)=K^{\mathsf t}B(PY^{\mathsf t}+YP^{\mathsf t})K$で計算できる。
被覆の接空間$\ker G$上でこれは交代行列全体に全射である。
具体的には、前記検算で$T+T^{\mathsf t}=0$、上三角六成分の線型階数6、
$\operatorname{Pf}(T)=15q/16$となる（定数は検算の基底による）。
滑らかな被覆の垂直方向は$T$で消えるので、元の接空間でもその階数は6である。

系\ref{cor:reversible-quadratic-cone}の局所表示により、$I$は六つの法線座標で生成され、
その一次部分は独立である。$\mathfrak m$を極大イデアルとすれば、$S$の
$I/\mathfrak m I$での成分は上記の一般交代行列の六座標となる。中山の補題から$S$の成分は$I$を生成する。
一般の四行交代行列の二次小行列式は六座標の全二次式を張る
（係数行列の階数は21）。従って$S$の二次小行列式の像は
$I^2/\mathfrak m I^2$全体を張り、再び中山の補題でそのイデアルは$I^2$になる。
六行を消去した後のこれら二つのイデアルが第三階数の各小行列式イデアルそのものである。
\end{proof}
